\documentclass[12pt]{article}
\usepackage[T1]{fontenc}
\usepackage[utf8]{inputenc}
\usepackage{lmodern}
\usepackage{textcomp}
\usepackage{enumitem}
\usepackage{geometry}
\geometry{margin=1in}
\begin{document}
\begin{center}
Answer Key - Astronomy - Graduate Level Assessment 4 \
\small (Answers are ordered exactly as in the question paper.)
\end{center}

\section*{Multiple Choice}
\begin{enumerate}[label=\arabic*.]
\item \textbf{Answer:} A
\\ \textit{Options:} A) asteroids would orbit with a period half that of Jupiter; B) asteroids would orbit with a period twice that of Jupiter; C) asteroids would orbit with a period twice that of Mars; D) A and B
\\ \textit{Note:} Kirkwood gaps occur in the asteroid belt at positions where the gravitational influence of Jupiter creates resonances that destabilize the orbits of asteroids, specifically at locations where their orbital period is a simple fraction, such as half, of Jupiter's orbital period.
\item \textbf{Answer:} D
\\ \textit{Options:} A) fast rotation only; B) a rocky mantle only; C) a molten metallic core only; D) both a molten metallic core and reasonably fast rotation
\\ \textit{Note:} A terrestrial planet requires a molten metallic core to facilitate the movement of electrically conductive fluids, which, combined with reasonably fast rotation, generates a dynamo effect that produces a strong magnetic field.
\item \textbf{Answer:} C
\\ \textit{Options:} A) a New Horizons flyby in the 1990 s; B) Hubble Space Telescope images that resolved Pluto's disk; C) brightness measurements made during mutual eclipses of Pluto and Charon; D) radar observations made by the Arecibo telescope
\\ \textit{Note:} The correct answer is based on the fact that brightness measurements during mutual eclipses of Pluto and its moon Charon allowed astronomers to determine the size of Pluto by observing the changes in luminosity as the two bodies obscured each other, providing a direct method to calculate its diameter.
\item \textbf{Answer:} C
\\ \textit{Options:} A) a composition dominated by hydrogen and helium; B) the presence of belts zones and storms; C) an equatorial wind speed of more than 900 miles per hour; D) significant "shear" between bands of circulation at different latitudes
\\ \textit{Note:} The equatorial wind speed of more than 900 miles per hour is not a similarity between Saturn and Jupiter's atmospheres because while Jupiter is known for its high wind speeds, Saturn's equatorial winds are significantly slower, often averaging around 1,100 kilometers per hour (about 680 miles per hour).
\item \textbf{Answer:} C
\\ \textit{Options:} A) the entire planets are made mostly of metal.; B) metals condensed first in the solar nebula and the rocks then accreted around them.; C) metals sank to the center during a time when the interiors were molten throughout.; D) radioactivity created metals in the core from the decay of uranium.
\\ \textit{Note:} The correct answer is based on the principle of density differentiation, where heavier metals, such as iron and nickel, migrated to the center of the terrestrial planets during their formation when their interiors were still molten, allowing for the stratification of materials according to density.
\end{enumerate}

\section*{True / False}
\begin{enumerate}[label=\arabic*.]
\setcounter{enumi}{5}
\item \textbf{Answer:} False
\\ \textit{Options:} A) True; B) False
\\ \textit{Note:} The solar cycle typically spans approximately 11 years, not 9, resulting in a peak and a subsequent decline in sunspot activity over that period.
\item \textbf{Answer:} False
\\ \textit{Options:} A) True; B) False
\\ \textit{Note:} The lithosphere is primarily composed of the rigid, solid rock of the Earth's crust, not the softer, more ductile materials found in the underlying mantle.
\item \textbf{Answer:} True
\\ \textit{Options:} A) True; B) False
\\ \textit{Note:} The statement is true because the cool hydrogen gas in the star's atmosphere absorbs specific wavelengths of light from the star's blackbody spectrum, resulting in absorption lines characteristic of hydrogen in the visible spectrum.
\item \textbf{Answer:} False
\\ \textit{Options:} A) True; B) False
\\ \textit{Note:} Impact heating is not a significant mechanism for maintaining a stable atmosphere on Mercury because the planet's weak gravitational field and extreme temperatures lead to a lack of a substantial, lasting atmosphere, regardless of solar proximity or impact events.
\item \textbf{Answer:} False
\\ \textit{Options:} A) True; B) False
\\ \textit{Note:} Iron is not the second most common element in the solar system; instead, helium ranks second after hydrogen in terms of abundance.
\end{enumerate}

\section*{Long Form Response}
\begin{enumerate}[label=\arabic*.]
\setcounter{enumi}{10}
\item \textbf{Answer:} The relationship between mass, density, and size in gas giants is fundamentally influenced by gravitational compression. Jupiter, being the most massive planet in our solar system, experiences significant gravitational forces that compress its gaseous envelope, resulting in a higher density compared to Saturn. Despite Saturn's larger size, its lower mass means it experiences less gravitational compression, leading to a less dense structure. Consequently, while Saturn's volume is nearly comparable to Jupiter's, its mass is insufficient to induce the same degree of compression, illustrating how mass directly impacts the density and overall size of gas giants. This relationship highlights the unique balance between gravitational forces and the physical properties of gaseous planets.
\\ \textit{Note:} The answer correctly identifies that the difference in mass between Jupiter and Saturn leads to varying levels of gravitational compression, affecting their densities and ultimately explaining why Saturn, despite having a lower mass, can maintain a size nearly comparable to that of Jupiter.
\item \textbf{Answer:} Electromagnetic radiation is arranged from shortest to longest wavelength as follows: gamma rays, X-rays, ultraviolet, visible light, infrared, and radio waves. Each category plays a crucial role in astronomical observations and applications. Gamma rays and X-rays are vital for studying high-energy phenomena such as black holes and supernovae, while ultraviolet light reveals information about the formation and evolution of stars. Visible light allows for direct observations of celestial bodies, whereas infrared radiation helps astronomers investigate cooler objects, such as dust clouds and distant galaxies. Finally, radio waves are essential for mapping cosmic structures and understanding the universe's large-scale properties, making each segment of the spectrum integral to comprehensive astronomical research.
\\ \textit{Note:} correct because it accurately lists the categories of electromagnetic radiation in order of wavelength and emphasizes their significance in astronomical observations, detailing how each category contributes to understanding various cosmic phenomena.
\end{enumerate}

\vfill
\noindent\textit{End of Answer Key}
\end{document}